\subsection*{一步一步建立直观}

\paragraph{1. 先把对数读作倍率的坐标}

对数回答“要使用多少指数参数，才能得到这个正倍率”。乘法在原尺度上组合倍率，在对数尺度上则变成加法。右图选择自然对数，用曲线展示这种坐标压缩：相同的绝对输入增量，作用在较大基数上时代表较小的相对变化。



\paragraph{2. 用其他函数作对照，辨认曲率}

切换到平方函数，看到输入越大，局部变化率反而越大。这与对数曲线越来越平缓形成对照。函数图像的弯曲不是装饰，它反映输出更关心绝对变化还是相对变化。不过乘法变加法的恒等式仍需由指数与对数定义推导，不能只凭曲线形状断言。



\paragraph{3. 定义域是意义的一部分}

回到对数函数并移动观察点。实对数要求输入为正，因为它在这里是正倍率的指数坐标；接近零时导数变大，非正位置则没有实函数值。复数对数可以扩展讨论，却需要处理辐角分支，不能将它与实对数公式不加说明地混用。


